\paragraph{} Este apêndice dedica-se a estimar o tempo necessário para a obtenção de uma solução ótima, por força bruta, referente à instância RC201, a qual contém 100 clientes (incluindo o depósito). O objetivo fundamental desta seção é evidenciar a natureza combinatória da complexidade temporal intrínseca a problemas de otimização de rotas \textit{NP-hard} \cite{Garey79}, como o \textit{VRP}. Mesmo adotando a simplificação extrema do cenário para uma formulação análoga ao clássico Problema do Caixeiro Viajante (\textit{TSP}), onde um único veículo visita sequencialmente todos os destinos, as permutações resultantes exigem um esforço computacional assintoticamente proibitivo.

\paragraph{} O número de sequências possíveis para um veículo visitar todos os pontos e retornar à origem é dado por:
\[
    99! \approx 9,3326 \times 10^{155}
\]

\paragraph{} Para este cenário, assumimos um supercomputador teórico capaz de avaliar 1 bilhão ($10^9$) de soluções por milissegundo. Calculando a vazão anual desta máquina:
\[
    10^{12} \times 60 \times 60 \times 24 \times 365,25 \approx 3,15 \times 10^{19} \text{ sol./ano}
\]

\paragraph{} O tempo estimado para encontrar o ótimo absoluto seria:
\[
    \text{Tempo} = \frac{9,33 \times 10^{155}}{3,15 \times 10^{19}} \approx 2,96 \times 10^{136} \text{ anos}
\]

\paragraph{} Para compreender a magnitude deste resultado, a idade estimada do universo é de aproximadamente $13,8 \times 10^9$ anos. Outra comparação é que há ``apenas'' $10^{80}$ átomos no universo observável. Se tivéssemos começado o processamento no momento do \textit{Big Bang}, até o presente momento teríamos percorrido uma fração tão pequena do espaço de busca que ela seria estatisticamente equivalente a $0\%$.